% rig_results.tex -- Results: hardware (draft 4 Oct 2026). No run ids in the text.
% Sources: docs/rig_2026-09-30_review.md, rig_2026-09-30_replay.md, rig_2026-09-30_force_balance.md,
% registry RIG-0930-* and RIG-1001-* (h1, h2-z080, c1-z100, f8-z100, h5-z100-rot120, h6-z100-xp60),
% logs results/rig/2026-10-01/<flight>_logs; B4 offline replay (2026-10-01_load_ext_force.md);
% CURRENT_STATE 7 items 8 and 10 (16 and 23 Sep); twin R0817/R0818, R0834-R0839.
% The H6 drift is 2.8 cm in the registry and 4.5 cm in the 3 Oct review recompute: reported as 3-5.
% FIGURE TODO (thesis-figure skill): 1 Oct ring height vs time for H1/H2/C1/H5 with the cap and the
% integral; tilt vs ring height; plan view of the -x drift with the rotation test.

\subsection{Hardware results}
\label{sec:rig_results}

The hardware campaign so far covers carrying the ring with three and four UAVs. Neither detach nor
attach has been flown on the rig. The rig is available about two days a week, and each day
followed the simulation work of the days before it, so each is summarised with what it showed and
what changed as a result.

\subsubsection{Rig days}

\begin{table}[h]
\centering
\small
\caption{Rig sessions with the payload.}
\label{tab:rig_days}
\begin{tabularx}{\textwidth}{lX}
\toprule
Date & What was flown and what it showed \\
\midrule
3 Aug & First flights of the planner stack on the rig. Fixed: attachment slots matched in the
payload's frame (a yawed payload made the problem infeasible), a latched payload yaw datum, a latched
heading per UAV (one UAV turned a full circle on takeoff), trackers on the correct clock. \\
16 Sep & First carry attempts with the planner from the floor: the payload did not lift. One UAV per
flight spun in yaw at spool-up, later traced to the tracker's yaw command on a tilted floor rest. The planner's references assumed taut rods and sat inside
the UAVs (Section~\ref{sec:takeoff}). \\
23 Sep & Two UAVs, rods typed 0.50\,m against a real 0.46\,m, thrust gain typed: the payload did not
leave the floor and the UAVs drifted in over it. One battery sagged from 21.2 to 17\,V. Found: battery
thresholds set for a smaller battery, and the need for the rod sweep. \\
30 Sep & Four UAVs, 0.86\,kg ring. First lifts and the first sustained carry (12.6\,s at 0.28\,m). Seven
bugs fixed on the day. The thrust map identified (Section~\ref{sec:hardware}). Breakaway at 60--80\,\%
of the planned pull. Last flight: ring held 16\,cm low with a 14\,cm heave at 0.17\,Hz, and 25 planner
solve failures. \\
1 Oct & Three UAVs (one airframe out of service), six flights to 1.0\,m including a circle and a
figure-eight, with no solve failures (Table~\ref{tab:rig_oct1}). \\
\bottomrule
\end{tabularx}
\end{table}

\paragraph{30 September.} The day went from a ring that would not leave the floor to sustained
carries. The bugs found on the day each had a clear signature in the logs. The creep fed the UAVs to the
wrong plates, so the fleet flew into the middle of the ring. The cable feedforward was zeroed while the
ring rested, so the lift deadlocked. Only one UAV received the pull flag. The pull arrived as a step
and launched the ring, releasing two magnets. Each was fixed and then flown, which is how the
pretension and breakaway procedure of Section~\ref{sec:takeoff} came about. Two problems remained at
the end of the day: the ring was held low, and the planner failed to solve 25 times in the last
flight. A replay of the planner on the logged states traced the failures to two faults acting
together. The UAVs sat 30$^\circ$ round from the attachment points the planner modelled, one plate off,
and the payload's measured quaternion changed sign when its yaw crossed 180$^\circ$. Removing either
removed every failure in the replay. The planner now keeps the quaternion's sign continuous and
warns before ARM if a UAV sits more than 10$^\circ$ from its slot. None of the 1 October flights had a
solve failure.

\paragraph{Thrust identification.} Each UAV hovered free with masses of 0, 98 and 196\,g hung from it,
on batteries between 22.7 and 24.6\,V. The 13 points fit the affine map of
Section~\ref{sec:hardware}, $u_T = a_i + 0.507 M - 0.022 (V - 23.5)$, with offsets $a_i$ of 0.181 to
0.187 and a residual of 0.003 rms. With the map and one shared offset, the four UAVs hovered free at
+3.7, $-0.5$, $-0.4$ and +0.9\,cm from their targets with no integral action. In tethered flight the
throttle the UAVs needed (0.55--0.60) matched what the map predicted for their share of the load
(0.569--0.575).

\subsubsection{1 October: three UAVs to 1.0\,m}

\begin{table}[h]
\centering
\small
\caption{Rig flights on 1 October, three UAVs on plates 0/4/8. Breakaway: fraction of the planned
pull at which the ring left the floor; it was held for the whole flight. Height integral bound
0.15\,m. Tilt and the horizontal offset are over the hold; $^{\dagger}$ marks the mean distance from
the trajectory reference over the manoeuvre. Orbit: radius 0.5\,m at 0.125\,m/s; figure-eight:
$a = 0.5$\,m at 0.2\,m/s. Rotated: the whole system turned by 124$^\circ$; moved: shifted 1.1\,m
towards $+x$.}
\label{tab:rig_oct1}
\footnotesize
\begin{tabularx}{\textwidth}{>{\raggedright\arraybackslash}Xcccccc}
\toprule
Flight & Target & Held & Integral & Tilt & Offset & Break- \\
 & (m) & (m) & & mean / max & (cm) & away \\
\midrule
Hover & 0.5 & 0.45 & off & 5.0 / 13.8$^\circ$ & 5.7 & 80\,\% \\
Hover & 0.8 & 0.67 & at bound & 2.5 / 6.7$^\circ$ & 11.8 & 73\,\% \\
Orbit & 1.0 & 0.94 & at bound & 1.6 / 4.1$^\circ$ & 15.1$^{\dagger}$ & 83\,\% \\
Figure-eight & 1.0 & 0.67 & never on & 3.2 / 9.0$^\circ$ & 18.8$^{\dagger}$ & 73\,\% \\
Hover, rotated & 1.0 & 1.01 & at bound & 3.8 / 6.6$^\circ$ & 8.1 & 93\,\% \\
Hover, moved & 1.0 & 0.89 & at bound & 2.2 / 4.1$^\circ$ & 3--5 & 77\,\% \\
\bottomrule
\end{tabularx}
\end{table}

All six flights lifted cleanly. The rod sweep handed over at 51--55$^\circ$, as in the twin, and the
measured rods read 0.531--0.556\,m against 0.55\,m typed. The orbit completed a full lap and the
figure-eight its whole path, and every flight landed. Three results stand out.

\paragraph{The ring sat low.} Without the height integral the ring held 5\,cm low at 0.5\,m, and at
higher targets it sat 20--30\,cm low before the integral took over, which then saturated at its bound.
The figure-eight, where the integral does not run, held the ring 29\,cm low. The shortfall was largest
on the flights whose ring broke free at the smallest fraction of the planned pull. The cause was
found afterwards in the logs: the breakaway freeze had held the pull at that fraction for the whole
flight, so the trackers were told the rods pull 7--27\,\% less than they do. Holding the pull at
75\,\% in the Gazebo twin reproduced both symptoms on the rig: the ring 23--26\,cm low and the rods
steepened from 52 to 58$^\circ$. Releasing the cap once the lift starts restored the twin's height
to the no-cap value (Section~\ref{sec:sim_to_real}). This fix has not yet been flown.

\paragraph{Tilt fell with height.} Over the first three flights the ring's mean tilt fell from
5.0$^\circ$ at 0.45\,m to 2.5$^\circ$ at 0.67\,m and 1.6$^\circ$ at 0.94\,m; the later hovers at
0.9--1.0\,m had 2.2--3.8$^\circ$. Close to the floor the ring sits in the UAVs' downwash and its
recirculation off the floor, and flying higher was the most effective single change for a steady
ring.

\paragraph{The ring drifted towards $-x$.} In every flight the ring sat 3--17\,cm from its horizontal
reference, always towards the room's $-x$ direction, and nothing in the controller pulled it back.
The offset kept its direction when the whole system (ring and UAVs) was rotated by 124$^\circ$, so it
does not come from a particular UAV, rod or the ring. The force behind it was estimated offline from
the logs: the UAVs' delivered thrust, from the identified map and their measured attitudes, less the
total weight. Over the steady holds this gives a horizontal force of 0.41--0.52\,N towards $-x$,
about 1$^\circ$ of the fleet's thrust. At the planner's horizontal stiffness of about 5\,N/m that
explains offsets of 8--10\,cm. Moving the system 1.1\,m towards $+x$ reduced the offset to 3--5\,cm,
in a single flight. Two explanations remain open: a ground plane in the motion-capture calibration
tilted by about 1$^\circ$, which would put a component of gravity along $x$ in the frame the
controller works in, or air moving in the room. A plumb-line check of the calibration comes first at
the next session.

\subsubsection{Simulation against the rig}
\label{sec:sim_to_real}

Table~\ref{tab:sim_vs_rig} sets the twin's hover at 1.0\,m against the rig's flights at the same
target. The twin comparison of the breakaway cap was flown with four UAVs and the rig with three, so
the comparison is of mechanisms rather than of exact values. The twin matched the rig in the takeoff
geometry and in solver health, and a single setting, the breakaway cap, accounted for the height
difference. The twin did not show the horizontal drift, the
larger tilt at low height or the slower breakaway.

\begin{table}[h]
\centering
\small
\caption{Hover at 1.0\,m: Gazebo rig twin (four UAVs, before the tracker fix) against the rig on
1 October (three UAVs). The twin with the cap held at 0.75 reproduces the rig's pull.}
\label{tab:sim_vs_rig}
\begin{tabularx}{\textwidth}{>{\raggedright\arraybackslash}Xccc}
\toprule
 & Twin, cap released & Twin, cap held at 0.75 & Rig, 1 October \\
\midrule
hand-over elevation & 55.4--55.7$^\circ$ & 55--56$^\circ$ & 51--55$^\circ$ \\
measured rods (m) & 0.550 & 0.550 & 0.531--0.556 \\
breakaway, fraction of the pull & about 0.97 & 0.75 (set) & 0.73--0.93 \\
ring height before the integral (m) & 1.03 & 0.77 & 0.71--0.81 \\
rod elevation in the hold & 52$^\circ$ & 58$^\circ$ & 57--65$^\circ$ \\
ring tilt, mean & 1.2$^\circ$ & 1.3$^\circ$ & 1.6--3.8$^\circ$ \\
horizontal offset in hover & under 1\,cm & under 1\,cm & 3--17\,cm \\
planner solve failures & 0 & 0 & 0 \\
\bottomrule
\end{tabularx}
\end{table}
